\section{Evaluated Component Matrix}

We do not introduce a new architecture. Instead, we define a controlled component matrix that isolates the effect of SupCon, ArcFace, BNNeck, and their combination under identical backbone, optimizer, embedding dimensionality, checkpoint-selection, and evaluation protocols.

\subsection{Component Matrix and Evaluation Scope}

This work evaluates a controlled family of ResNet18 embedding models rather than proposing a new palmprint architecture. All tested variants use the same backbone family and a 256-D embedding size; they differ only in the use of supervised contrastive learning, BNNeck, ArcFace, and a light SupCon term. The main variants are:
\begin{itemize}
    \item B0: ResNet18 + cross entropy;
    \item B1: ResNet18 + cross entropy + supervised contrastive learning;
    \item B4: ResNet18 + ArcFace;
    \item B5: ResNet18 + BNNeck + cross entropy;
    \item B6: ResNet18 + BNNeck + ArcFace;
    \item B7: ResNet18 + BNNeck + ArcFace + light supervised contrastive learning.
\end{itemize}
B1 is the current CE + SupCon baseline. B6 is the initially hypothesized BNNeck + ArcFace variant because it appeared favorable in the seen-identity diagnostic setting. B5 is included to isolate the effect of BNNeck without ArcFace. The strict Tongji component matrix is used to distinguish the original B6 hypothesis from the broader component behavior observed under palm-class-disjoint evaluation.

\subsection{BNNeck Embedding}

BNNeck was proposed as a normalization neck for separating the feature space used for metric comparison from the feature space optimized by classification supervision \cite{luo2019strong}. Let $h_i$ denote the output feature of the ResNet18 backbone for sample $i$. A linear projection maps this feature to a 256-D vector:
\begin{equation}
    z_i = W_e h_i + b_e.
\end{equation}
BNNeck applies batch normalization to this projected vector:
\begin{equation}
    \tilde{z}_i = \mathrm{BN}(z_i).
\end{equation}
For BNNeck variants, the evaluation embedding is the post-BN representation:
\begin{equation}
    f_i = \frac{\tilde{z}_i}{\lVert \tilde{z}_i \rVert_2}.
\end{equation}
For non-BNNeck variants, the corresponding 256-D embedding is L2-normalized before matching. Gallery and probe embeddings are compared with cosine similarity.

\subsection{ArcFace Supervision}

ArcFace uses normalized features and normalized class weights to impose an additive angular margin \cite{deng2019arcface}. Let $W_j$ be the normalized classifier weight for class $j$, and let $\theta_j$ be the angle between the normalized embedding $f_i$ and $W_j$. For the ground-truth class $y_i$, ArcFace adds an angular margin $m$ and scales logits by $s$:
\begin{equation}
    \ell_i =
    -\log
    \frac{
        \exp \left(s \cos(\theta_{y_i} + m)\right)
    }{
        \exp \left(s \cos(\theta_{y_i} + m)\right)
        + \sum_{j \ne y_i} \exp \left(s \cos(\theta_j)\right)
    }.
\end{equation}
In the current B6 configuration, $s=30.0$ and $m=0.5$. ArcFace directly shapes angular classifier geometry, but the experiments test whether that training geometry actually improves held-out palm-class verification and identification metrics.

\subsection{What Is Not Claimed as the Main Method}

Earlier Gabor variants remain useful baselines and ablations because Gabor-style orientation coding is historically important in palmprint recognition \cite{kong2004competitive,genovese2019palmnet,liang2021compnet}, but they are not the main supported claim in this paper. The present evidence does not support a fixed-Gabor, learnable-Gabor, universal BNNeck + ArcFace, state-of-the-art, or cross-dataset robustness claim. The supported claim is a scoped component evaluation under the stated Tongji and IITD protocols.
